\section*{Abstract}
	Das T0-Modell interpretiert die kosmologische Rotverschiebung nicht als Expansion, sondern als Energieverlust von Photonen im universellen $\xi$-Feld. Dieses Dokument behandelt den feldtheoretischen Energieverlust-Mechanismus und die Verbindung zur Hubble-Konstante. Die Skala der Rotverschiebung ist über den gemessenen $H_0$ kalibriert (P39); als Strukturform trägt der Korpus $\hbar H_0/(m_ec^2) = \frac{\pi}{2}\xi^{10}$ mit $H_0 \approx 66{,}82\,$km/s/Mpc (Dok.~165, 388). Die Einheitenumrechnung, die Frequenzunabhängigkeit und die Hubble-Spannung sind in Dokument~026 (Geometrische Kosmologie) behandelt.


\section{Einleitung: Die Hubble-Konstante neu gedacht}

	Die konventionelle Interpretation des Hubble-Gesetzes geht davon aus, dass sich Galaxien aufgrund des expandierenden Raums voneinander entfernen ($v = H_0 d$). Dieses Expansionsparadigma erfordert 69\% dunkle Energie, leidet unter anhaltenden Messspannungen und Feinabstimmungsproblemen.

	Das T0-Modell bietet eine fundamental andere Perspektive: Das Universum ist statisch, und die beobachtete Rotverschiebung entsteht durch Energieverlust von Photonen während ihrer Ausbreitung durch das universelle $\xi$-Feld. Die Hubble-Konstante wird damit von einem Maß für Raumexpansion zu einer charakteristischen Energieverlustrate.

	Die Zeit-Energie-Dualität ($\Delta E \cdot \Delta t \geq \hbar/2$) verbietet einen zeitlichen Beginn des Universums: Ein endliches Zeitintervall ab einer Singularität würde unendliche Energieunschärfe erfordern. Das Universum muss ewig existiert haben, wodurch Raumexpansion zur Erklärung kosmischer Beobachtungen überflüssig wird.

\section{Das universelle $\xi$-Feld}

	Der Grundstein des T0-Modells ist die universelle geometrische Konstante:
	\begin{equation}
		\xi = \frac{4}{3} \times 10^{-4} = 1{,}3333\ldots \times 10^{-4}
	\end{equation}

	Diese dimensionslose Konstante bestimmt eine charakteristische Energieskala:
	\begin{equation}
		E_\xi = \frac{1}{\xi} = 7500 \quad \text{(natürliche Einheiten)}
	\end{equation}

	Das $\xi$-Feld durchdringt den gesamten Raum und vermittelt Wechselwirkungen zwischen Photonen und dem Vakuum. Seine Dynamik wird durch die Wellengleichung beschrieben:
	\begin{equation}
		\square E_{\text{field}} = \left(\nabla^2 - \frac{\partial^2}{\partial t^2}\right) E_{\text{field}} = 0
	\end{equation}

\section{Energieverlust-Mechanismus}

	\subsection{Fundamentale Energieverlust-Gleichung}

	Photonen verlieren Energie durch Wechselwirkung mit dem $\xi$-Feld. Die Rate des Energieverlusts folgt der Differentialgleichung:
	\begin{equation}
		\frac{dE}{dx} = -\frac{H_0}{c} \cdot E
	\end{equation}

	Das negative Vorzeichen zeigt Energieverlust an; die lineare Abhängigkeit von $E$ stellt sicher, dass die resultierende Rotverschiebung frequenzunabhängig ist (siehe Dokument~026 für die geometrische Begründung). Der Vorfaktor muss eine inverse Länge sein: Mit dem dimensionslosen $\xi$ allein hätte $-\xi E$ die Dimension $[E]$, $dE/dx$ in natürlichen Einheiten aber $[E^2]$; ebenso hätte $\xi \cdot x$ die Dimension einer Länge ($[E^{-1}]$). Mit $x$ in Metern ergäbe $z = \xi x$ bereits bei $x = 1/\xi = 7500$~m den Wert $z = 1$. Die Gleichungen werden deshalb mit der inversen Länge $H_0/c$ geschrieben.

	\subsection{Lösung für kosmologische Entfernungen}

	Für kosmologische Beobachtungen ($H_0 x/c \ll 1$) ergibt die störungstheoretische Lösung:
	\begin{equation}
		E(x) = E_0 \left(1 - \frac{H_0}{c} \cdot x\right)
	\end{equation}

	Die Rotverschiebung ist damit:
	\begin{equation}
		z = \frac{E_0 - E(x)}{E(x)} \approx \frac{H_0\, x}{c}
	\end{equation}

	Diese Beziehung reproduziert das beobachtete lineare Hubble-Gesetz ohne Raumexpansion.

\section{Verbindung zur Hubble-Konstante}

	Der Vergleich mit dem Hubble-Gesetz $z = H_0 d/c$ legt die Skala des Energieverlusts nur über den gemessenen $H_0$ fest; $H_0$ ist in diesem Sinn kalibriert, nicht aus $\xi$ allein hergeleitet (P39). Als Strukturrelation trägt der Korpus (Dok.~165):
	\begin{equation}
		\boxed{H_0^{\text{T0}} = \frac{\pi}{2}\,\xi^{10}\,\frac{m_e c^2}{\hbar}}
	\end{equation}

	\textbf{Ergebnis:} $H_0^{\text{T0}} \approx 66{,}82\,$km/s/Mpc ($-0{,}86\%$ gegenüber Planck). Exponent~10 und Vorfaktor $\pi/2$ sind nicht hergeleitet.

	Die ältere Form $H_0 = E_H/\hbar$ mit $E_H = E_0\,\xi^{41/4}$ ($\approx 66{,}2\,$km/s/Mpc) gilt nicht: Über die in Dok.~388 als Kandidat geführte Beziehung $T_{\mathrm{CMB}}^4 = 16\,H_0\,m_e^3$ ergibt sie $T_{\mathrm{CMB}} = 2{,}71898$~K, $11\sigma$ neben dem FIRAS-Wert, während die Form mit Exponent~10 auf $0{,}1\sigma$ trifft.

	Für die schrittweise Einheitenumrechnung (natürliche Einheiten $\rightarrow$ SI $\rightarrow$ km/s/Mpc), den Vergleich mit experimentellen Messwerten und die Hubble-Spannung siehe \textbf{Dokument~026} (Geometrische Kosmologie, Abschnitte 4--6).

\section{Dimensionsanalyse}

	Verifikation der Dimensionskonsistenz in natürlichen Einheiten:

	\begin{center}
		\adjustbox{max width=\textwidth}{%
\begin{tabular}{lccc}
			\toprule
			\textbf{Größe} & \textbf{T0-Ausdruck} & \textbf{Dimension} & \textbf{Status} \\
			\midrule
			Geometrische Konstante & $\xi = 4/3 \times 10^{-4}$ & $[1]$ & \checkmark \\
			Energieskala & $E_\xi = 1/\xi$ & $[E]$ & \checkmark \\
			Energieverlustrate & $dE/dx = -(H_0/c) \cdot E$ & $[E^2]$ & \checkmark \\
			Rotverschiebung & $z = H_0 x/c$ & $[1]$ & \checkmark \\
			Hubble-Parameter & $H_0 = \frac{\pi}{2}\xi^{10} m_ec^2/\hbar$ & $[T^{-1}]$ & \checkmark \\
			\bottomrule
		\end{tabular}%
}
	\end{center}
	Die Zeilen Energieverlustrate und Rotverschiebung sind nur mit der inversen Länge $H_0/c$ konsistent; die Form mit $\xi$ allein ($-\xi \cdot E$ mit Dimension $[E]$, $\xi \cdot x$ mit Dimension $[E^{-1}]$) gilt nicht (siehe Abschnitt Energieverlust-Mechanismus).

\section{Theoretische Vorteile}

	Die Neuinterpretation der Hubble-Konstante als Energieverlustrate löst mehrere Probleme:

	\begin{enumerate}
		\item \textbf{Eliminierung dunkler Energie:} Die scheinbare kosmische Beschleunigung entsteht natürlich aus dem entfernungsabhängigen Energieverlust. Keine exotische Energieform mit negativem Druck ist erforderlich.

		\item \textbf{Feinabstimmungsprobleme gelöst:} Das Horizontproblem verschwindet, da alle Bereiche über unendliche Zeit in kausalem Kontakt waren. Das Flachheitsproblem entfällt mangels anfänglicher Bedingungen.

		\item \textbf{Parameterreduktion:} Wo $\Lambda$CDM sechs unabhängige Parameter benötigt, führt T0 die kosmologischen Observablen auf die geometrische Konstante $\xi$ zurück; die Skala der Rotverschiebung ist dabei über den gemessenen $H_0$ festgelegt und folgt nicht aus $\xi$ allein.
	\end{enumerate}

\section{Fazit}

	\begin{equation}
		\boxed{H_0^{\text{T0}} = \frac{\pi}{2}\,\xi^{10}\,\frac{m_e c^2}{\hbar} \approx 66{,}82\,\text{km/s/Mpc}}
	\end{equation}

	Das Universum expandiert nicht. Die Hubble-Konstante misst Energieverlust im $\xi$-Feld, nicht Flucht; ihre Skala ist über den gemessenen Wert kalibriert (P39), die Strukturform mit Exponent~10 ist durch die CMB-Temperatur gestützt (Dok.~388). Die Finite-Elemente-Simulation des T0-Vakuums, die Frequenzunabhängigkeit als geometrischer Beweis und die Hubble-Spannung finden sich in Dokument~026.

	\begin{thebibliography}{99}

		\bibitem{pascher_geometrische_kosmologie}
		J. Pascher, \textit{T0-Kosmologie: Rotverschiebung als geometrischer Pfad-Effekt in einem statischen Universum}, T0-Dokumentenserie, Dok.~026.

		\bibitem{planck_2020}
		Planck Collaboration, \textit{Planck 2018 results. VI. Cosmological parameters}, Astron. Astrophys. 641, A6, 2020.

		\bibitem{riess_2022}
		A. G. Riess, et al., \textit{A Comprehensive Measurement of the Local Value of the Hubble Constant}, Astrophys. J. Lett. 934, L7, 2022.

	\end{thebibliography}
